\scanpage{99}\mspage{- 167 -}\hfill IX

$A$ \uline{lisse} sur $S$, et on désigne par
\[
 \underline E=\underline{\mathrm{Lie}}(A/S)\tag*{(13.1.1)}
\]
le module localement libre sur $S$, image inverse par la section nulle du faisceau tangent relatif de $A$ sur $S$. Nous admettrons ici la définition de l'homomorphisme exponentiel
\[
 \exp:E\longrightarrow A\quad,\tag*{(13.1.2)}
\]
où $E$ est le fibré vectoriel analytique sur $S$ dont le Module des sections holomorphes est $\underline E$. Cet homomorphisme se définit, au voisinage d'un point $s\in S$, d'abord au voisinage de la section nulle de $E$ par les méthodes bien connues en théorie des groupes formels en car. nulle, puis se prolonge de façon évidente par multiplicativité à $E$ tout entier. C'est un homomorphisme de groupes analytiques relatifs sur $S$, qui est \uline{étale} et a comme image le sous-espace ouvert $A^{\unsure{o|\circ}}$ de $A$ réunion des composantes neutres des fibres de $A$ sur $S$. Si $R$ est le sous-espace analytique noyau de $\exp$, on a donc une suite exacte de groupes analytiques relatifs :
\[
 0\longrightarrow R\longrightarrow E\longrightarrow A^o\longrightarrow0\quad,\tag*{(13.1.3)}
\]
où le morphisme $E\longrightarrow A^o$ est surjectif et étale, donc $R$ est un groupe analytique relatif \uline{étale} sur $S$. C'est un sous-espace analytique de $E$, qui est fermé dans $E$ si et seulement si $A^o$ est séparé sur $S$, i.e. la section nulle de $A^o$ est fermée dans $A^o$.
% archive sheet number: 98
% stamp: numbered label 485; 485 also in pencil above the lower strip.
